% ============================================================================
% ARBEITSBLATT: Kräfte zusammensetzen (UE 2)
% Physik Klasse 7 – Gesamtschule MV – 45 Minuten
% Version: Kariertes Papier mit TikZ, X-Kreuze
% ============================================================================
\documentclass[a4paper,11pt]{article}

\usepackage[left=1.2cm,right=1.2cm,top=1.5cm,bottom=1.5cm]{geometry}
\usepackage[utf8]{inputenc}
\usepackage[T1]{fontenc}
\usepackage[ngerman]{babel}
\usepackage{tikz}
\usepackage{fancyhdr}
\usepackage{tcolorbox}
\usepackage{amssymb,amsmath}
\usepackage{xcolor}

\renewcommand{\familydefault}{\sfdefault}

\pagestyle{fancy}
\fancyhf{}
\renewcommand{\headrulewidth}{0.4pt}
\lhead{\textbf{Physik Klasse 7}}
\chead{Kräfte zusammensetzen}
\rhead{Name: \underline{\hspace{3cm}}}
\setlength{\headheight}{14pt}
\cfoot{\thepage}

\definecolor{physik}{HTML}{2c5aa0}
\definecolor{gridcolor}{RGB}{200,200,200}
\definecolor{axiscolor}{RGB}{100,100,100}

\newcommand{\luecke}[1]{\underline{\hspace{#1}}}

% X-Kreuz (45° gedreht) am Angriffspunkt
% #1 = x-Position, #2 = y-Position
\newcommand{\xkreuz}[2]{%
    \draw[black, very thick] (#1-0.4,#2-0.4) -- (#1+0.4,#2+0.4);
    \draw[black, very thick] (#1-0.4,#2+0.4) -- (#1+0.4,#2-0.4);
}

\setlength{\parindent}{0pt}

\begin{document}

% ==================== MERKSATZ-BOX ====================

\begin{tcolorbox}[colback=white,colframe=physik,boxrule=1.5pt,arc=0pt,title={\textbf{Merksatz: Kräfte zusammensetzen}},left=3mm,right=3mm,top=2mm,bottom=2mm]

\textbf{Kraftpfeil:} \luecke{2.2cm}, \luecke{2.2cm}, \luecke{2.2cm}

\vspace{2mm}

\begin{minipage}[t]{0.48\textwidth}
\textbf{Gleiche Richtung:}\\[1mm]
$F_{\text{res}} = F_1 + F_2$
\end{minipage}
\hfill
\begin{minipage}[t]{0.48\textwidth}
\textbf{Entgegengesetzt:}\\[1mm]
$F_{\text{res}} = |F_1 - F_2|$
\end{minipage}

\vspace{2mm}

\textbf{Rechtwinklig:} $F_{\text{res}} = \sqrt{F_1^2 + F_2^2}$ \hfill \textbf{Maßstab:} 1 cm $\hat{=}$ 1 N

\end{tcolorbox}

\vspace{4mm}

% ==================== AUFGABE 1 ====================

\textbf{Aufgabe 1: Kraftpfeile zeichnen} \hfill (4 P)

Zeichne die Kraftpfeile im richtigen Maßstab (1 cm = 1 N = 2 Kästchen). Das $\times$ markiert den Angriffspunkt.

\vspace{3mm}

\begin{tabular}{@{}p{4.2cm}p{5.5cm}p{5.5cm}@{}}
a) $F = 3$ N nach rechts &
b) $F = 4$ N nach oben &
c) $F = 2{,}5$ N nach links \\[2mm]
\begin{tikzpicture}[scale=0.5]
    \draw[gridcolor, very thin] (0,0) grid (10,4);
    \xkreuz{2}{2}
\end{tikzpicture}
&
\begin{tikzpicture}[scale=0.5]
    \draw[gridcolor, very thin] (0,0) grid (10,10);
    \xkreuz{5}{1}
\end{tikzpicture}
&
\begin{tikzpicture}[scale=0.5]
    \draw[gridcolor, very thin] (0,0) grid (10,4);
    \xkreuz{8}{2}
\end{tikzpicture}
\end{tabular}

\vspace{6mm}

% ==================== AUFGABE 2 ====================

\textbf{Aufgabe 2: Kräfte in gleicher Richtung} \hfill (4 P)

Zwei Personen schieben einen Wagen: $F_1 = 3$ N $\rightarrow$, \quad $F_2 = 4$ N $\rightarrow$

\vspace{2mm}

a) Zeichne $F_1$ (blau), $F_2$ (grün) und $F_{\text{res}}$ (rot) ein: \hfill (2 P)

\vspace{2mm}

\begin{center}
\begin{tikzpicture}[scale=0.5]
    \draw[gridcolor, very thin] (0,0) grid (18,6);
    % Achse
    \draw[axiscolor, thick] (0,3) -- (18,3);
    % X-Kreuz am Angriffspunkt
    \xkreuz{2}{3}
    \node[below,font=\small] at (2,2.3) {Wagen};
    % Maßstab
    \node[font=\tiny,gray] at (9,0.5) {Maßstab: 2 Kästchen = 1 N};
\end{tikzpicture}
\end{center}

\vspace{2mm}

b) Berechnung: $F_{\text{res}} = $ \luecke{5cm} $= $ \luecke{2cm} \hfill (2 P)

\vspace{6mm}

% ==================== AUFGABE 3 ====================

\textbf{Aufgabe 3: Tauziehen} \hfill (5 P)

Team A zieht mit $F_1 = 6$ N nach links. Team B zieht mit $F_2 = 4$ N nach rechts.

\vspace{2mm}

a) Zeichne beide Kräfte und die Resultierende ein: \hfill (2 P)

\vspace{2mm}

\begin{center}
\begin{tikzpicture}[scale=0.5]
    \draw[gridcolor, very thin] (0,0) grid (22,6);
    % Achse
    \draw[axiscolor, thick] (0,3) -- (22,3);
    % X-Kreuz weiter rechts (Platz für 6N = 12 Kästchen nach links)
    \xkreuz{13}{3}
    % Teams
    \node[font=\small] at (3,5) {$\leftarrow$ Team A};
    \node[font=\small] at (20,5) {Team B $\rightarrow$};
\end{tikzpicture}
\end{center}

\vspace{2mm}

b) $F_{\text{res}} = $ \luecke{5cm} $= $ \luecke{2cm} \hfill (2 P)

c) Welches Team gewinnt? \luecke{3cm} \hfill (1 P)

\vspace{6mm}

% ==================== SEITENUMBRUCH ====================
\newpage

% ==================== AUFGABE 4 ====================

\textbf{Aufgabe 4: Kräftegleichgewicht} \hfill (2 P)

Beim Armdrücken drückt jeder mit $F = 5$ N.

\vspace{2mm}

a) $F_{\text{res}} = $ \luecke{4cm} \hfill (1 P)

\vspace{4mm}

b) Was passiert? \luecke{8cm} \hfill (1 P)

\vspace{8mm}

% ==================== AUFGABE 5 ====================

\textbf{Aufgabe 5: Rechtwinklige Kräfte} \hfill (3 P)

$F_1 = 3$ N nach rechts, \quad $F_2 = 4$ N nach oben

\vspace{2mm}

a) Zeichne beide Kräfte ein: \hfill (1 P)

\vspace{2mm}

\begin{center}
\begin{tikzpicture}[scale=0.5]
    \draw[gridcolor, very thin] (0,0) grid (12,12);
    % Achsen
    \draw[axiscolor, thick] (2,0) -- (2,12);
    \draw[axiscolor, thick] (0,2) -- (12,2);
    % X-Kreuz am Angriffspunkt
    \xkreuz{2}{2}
\end{tikzpicture}
\end{center}

\vspace{2mm}

b) Berechne mit Pythagoras: \hfill (2 P)

\vspace{1mm}

$F_{\text{res}} = \sqrt{F_1^2 + F_2^2} = \sqrt{3^2 + 4^2} = \sqrt{9 + 16} = \sqrt{25} = $ \luecke{1.5cm}

\end{document}
